\section{Parallel links and an endpoint-scale unrestricted error}
\label{sec:parallel-links-endpoint}

Fable's parallel-link lemma removes the forest excess from the link bound.
This section combines that lemma with two-scale restoration and a separate spectral estimate inside each clean block.
The endpoint parameters below give a relative error $O((\log r)^6/\sqrt r)$.
The argument is mathematical, with independent audits of its constituent estimates.

All logarithms in this section have base two unless written as $\ln$.

\begin{lemma}[Fable's parallel-link estimate]
\label{lem:endpoint-parallel-links}
Let $A,B$ be disjoint connected vertex sets in a connected multigraph $\Gamma$.
Let $P$ be a vertex set, and let $\xi$ be a uniform even subgraph.
Suppose its restriction outside $P$ is a forest with probability greater than one half.
For every integer $t\ge0$, at most $2^{t+1}$ components of $\Gamma-(A\cup B)$ satisfy all four conditions below:
\begin{enumerate}
\item They have an edge to each of $A,B$.
\item They avoid $P$.
\item They contain a cycle.
\item Their boundary has at most $t$ edges.
\end{enumerate}
\end{lemma}

\begin{proof}
If there is no component joining $A$ to $B$, the assertion is immediate.
Otherwise retain one such component $Q_0$.
The set $A\cup B\cup Q_0$ is connected.
For any two other complementary components $Q_i,Q_j$, the graph $\Gamma-Q_i-Q_j$ is connected.
It contains the retained connected set, and every other complementary component has an edge to that set.
No edge joins $Q_i$ and $Q_j$.

Thus the two regions have exactly one link.
The cycle-space pair law makes their zero-boundary events independent.
Conditional on both zero-boundary events, their internal even edge sets are independent and uniform.
This also follows from the direct product of the cycle spaces after deleting the two boundaries.

For each eligible component other than $Q_0$, let $E_i$ mean that its boundary is zero and its internal even edge set is nonzero.
The events $E_i$ are pairwise independent.
The zero-boundary condition imposes at most $t$ linear equations, so its probability is at least $2^{-t}$.
The cyclic internal graph has a nonzero even edge set with conditional probability at least one half.
Consequently $p_i:=\Pr(E_i)\ge2^{-t-1}$.

If there are $m$ such components, put $N=\sum_i\mathbf1_{E_i}$ and $\nu=\mathbb EN$.
Pairwise independence gives
\[
 \mathbb EN^2\le\nu^2+\nu,\qquad
 \Pr(N=0)\le\frac1{1+\nu},\qquad
 \nu\ge m2^{-t-1}.
\]
The forest event implies $N=0$.
Its probability exceeds one half, so $m<2^{t+1}$.
Including $Q_0$, if it was eligible, gives at most $2^{t+1}$ components.
\end{proof}

We recall the walk and moment estimates needed below.
Let $J$ have maximum degree three and girth greater than $D$.
Put $s=\lfloor(D-1)/2\rfloor$, and let $\mathcal C$ contain its cycles of length at most $L$.
For $a_C=2^{-|C|}$ and $\mu=\sum_Ca_C$, the local walk bound gives
\begin{equation}\label{eq:endpoint-local-weight}
 \sum_{C\ni v}a_C\le L2^{-s}.
\end{equation}
The repeated-walk estimate and the even spectral trace estimate give
\begin{equation}\label{eq:endpoint-trace-estimates}
 \mu\le\sum_{\ell\le L}\frac{\operatorname{tr}B_J^\ell}{2\ell2^\ell}
 \le(1+\tfrac32L^22^{-s})\mu,
 \qquad
 2^{-\ell}\operatorname{tr}B_J^\ell
 \ge\left(1-\frac{d_J}{2v_J}\right)^\ell-3v_J2^{-\ell/2}
\end{equation}
for even $\ell$, where $d_J=\sum_v(3-\deg_Jv)$ and $d_J<v_J$.
These are the established non-backtracking walk estimates from Section~\ref{sec:polynomial-foundations}.

Suppose $J=\Gamma-(P\cup S)$, where $S$ comes from a maximal disjoint packing of cycles of length at most $D$ outside $P$.
Write $c=\operatorname{comp}(\Gamma[P])$.
The usual link decomposition, with Lemma~\ref{lem:endpoint-parallel-links} replacing its cyclic-region term, gives
\begin{equation}\label{eq:endpoint-link-estimate}
 \ell(C,C')
 \le c+\frac{2L}{t+1}+2^{t+1}+578(L/D)^2
\end{equation}
whenever $8\le D\le L$ and $2(t+2\lfloor D/8\rfloor)\le D$.
Components meeting $P$ contribute at most $c$.
Components with more than $t$ boundary edges contribute at most $2L/(t+1)$.
The new lemma handles the other cyclic components.
Small tree components and direct edges are bounded by the unchanged girth argument.

\begin{lemma}[The refined second moment]
\label{lem:endpoint-two-link-moment}
Suppose every disjoint pair in $\mathcal C$ has at most $b$ links.
Put
\[
 \eta=\tfrac32\,2^bL^32^{-s}.
\]
Let $W_2$ be the total weight $a_Ca_{C'}$ of ordered disjoint pairs with exactly two links.
If the forest probability exceeds $1/2+\varepsilon$, then
\begin{equation}\label{eq:endpoint-moment-conclusions}
 W_2\ge\mu^2-(1+\eta)\mu,
 \qquad
 \mu\le\frac{1+\eta}{4\varepsilon}.
\end{equation}
The first inequality is strict when $\mu>0$.
\end{lemma}

\begin{proof}
If $\mu=0$, both conclusions hold immediately.
Assume $\mu>0$.
For each $C$, let $I_C$ mean that it is a component of $\xi$.
Put $p_C=\Pr(I_C)$ and $X_C=(a_C/p_C)\mathbf1_{I_C}$.
The cycle-component equations give $p_C\ge a_C$, so
\[
 \mathbb EX_C=a_C,\qquad \mathbb EX_C^2\le a_C.
\]
Distinct intersecting cycles have zero joint contribution.
For disjoint cycles, the pair law gives
\[
 \mathbb E(X_CX_{C'})=a_Ca_{C'}2^{\ell(C,C')-1}.
\]
The exceptional-partner bound and~\eqref{eq:endpoint-local-weight} make the total contribution from pairs with at least three links at most $\eta\mu$.
Thus, for $Z=\sum_CX_C$,
\[
 \mathbb EZ^2\le\mu^2+W_2+(1+\eta)\mu
 \le2\mu^2+(1+\eta)\mu.
\]
The forest event implies $Z=0$.
Cauchy--Schwarz gives
\[
 \mu^2\le\Pr(Z>0)\mathbb EZ^2
 <(1/2-\varepsilon)\mathbb EZ^2.
\]
The first displayed moment bound yields the strict lower bound on $W_2$.
The second yields the upper bound on $\mu$.
\end{proof}

\begin{lemma}[Cycle supply in one clean block]
\label{lem:endpoint-block-supply}
Fix $A=6e^{168}$.
Put $x=\log r$, $L=\lfloor Ax\rfloor$, and $s_0=\lfloor(D_0-1)/2\rfloor$.
Suppose $L^22^{-s_0}=o(1)$.
Let $J$ be a simple graph of maximum degree three, girth greater than $D_0$, and order $v$ satisfying
\[
 8L^2\le v\le3r,
 \qquad d_J\le2L.
\]
For sufficiently large $r$, its cycles of length at most $L$ have total weight at least one.
The threshold is uniform over all such graphs $J$.
\end{lemma}

\begin{proof}
For every even $\ell$ with $6x\le\ell\le L$, Bernoulli's inequality gives
\[
 \left(1-\frac{d_J}{2v}\right)^\ell
 \ge1-\frac{Ld_J}{2v}\ge\frac78.
\]
The nonreal spectral error in~\eqref{eq:endpoint-trace-estimates} is at most $9r^{-2}$.
The even harmonic sum over this interval tends to $\tfrac14\ln(A/6)$.
The repeated-walk factor tends uniformly to one.
Therefore the total cycle weight is at least
\[
 \left(\frac7{32}-o(1)\right)\ln(A/6)>1.
\]
No connectedness or eigenvalue multiplicity estimate is required.
\end{proof}

\begin{theorem}[An endpoint-scale forest error]
\label{thm:parallel-links-endpoint-forest}
Let $\Gamma$ be connected and subcubic, with rank $r$.
Let $P$ contain every vertex of degree less than three.
Put $x=\log r$ and $z=\log x$, and assume
\[
 |P|\le r/x^2,
 \qquad \operatorname{comp}(\Gamma[P])\le x/z^2.
\]
For a uniform even subgraph, let $\mathcal F$ be the event that its restriction to $\Gamma-P$ is a forest.
There is an absolute constant $M$ such that, for all sufficiently large $r$,
\begin{equation}\label{eq:endpoint-forest-log-six}
 \Pr(\mathcal F)\le\frac12+M\frac{x^6}{\sqrt r}.
\end{equation}
Separately, the more conservative parameter choice gives
\begin{equation}\label{eq:endpoint-forest-log-seven}
 \Pr(\mathcal F)\le\frac12+\frac{x^7}{\sqrt r}
\end{equation}
for all sufficiently large $r$.
\end{theorem}

\begin{proof}
If $P$ is empty, the forest event contains only the empty even edge set, so its probability is $2^{-r}$.
Assume $P$ is nonempty and fix one connected component $P_0$ as an anchor.
Keep $A=6e^{168}$, and put
\[
 L=\lfloor Ax\rfloor,
 \quad t=\lfloor z-2\log z\rfloor,
 \quad D_0=\left\lceil\frac{8Ax}{z}\right\rceil.
\]
We prove both bounds with the following two parameter rows:
\[
\begin{array}{c|c|c|c}
 &a&B&\varepsilon\\ \hline
 \text{conservative}&5&0&x^7/\sqrt r\\
 \text{endpoint}&4&\lceil\log(10^6A)\rceil&Mx^6/\sqrt r
\end{array}
\qquad
 D_1=\lfloor x-az-B\rfloor.
\]
The constant $M\ge1$ is chosen at the end, independently of $r$.
Suppose, for a contradiction, that $\Pr(\mathcal F)>1/2+\varepsilon$.
For large $r$, one has $8\le D_0<D_1\le L$ and $t\ge1$.

Choose a maximal disjoint family $\mathcal Q_0$ of cycles of length at most $D_0$ outside $P$.
Extend it to a maximal disjoint family $\mathcal Q_1$ for length $D_1$.
Let $S_j$ be their vertex unions and put $J_j=\Gamma-(P\cup S_j)$.
Then $S_0\subseteq S_1$, and $J_j$ is simple with girth greater than $D_j$.

Each chosen cycle is a cyclic region with boundary at most its length.
The exponential small-cut bound gives
\[
 |\mathcal Q_j|<8D_j2^{D_j}\ln(2/\varepsilon).
\]
Since $M\ge1$, both parameter rows have $\ln(2/\varepsilon)\le x$ for large $r$.
Consequently,
\begin{equation}\label{eq:endpoint-deletion-sizes}
 |S_0|\le8D_0^22^{D_0}x=r^{o(1)},\qquad
 |S_1|\le8\,2^{-B}\frac r{x^{a-3}}.
\end{equation}

Write $v_1=|V(J_1)|$ and $d_1=\sum_{v\in J_1}(3-\deg_{J_1}v)$.
The degree sum in $\Gamma$ gives $|V(\Gamma)|=2r+O(|P|)$, and hence $v_1=2r+o(r)$.
Also $d_1\le3(|P|+|S_1|)$.
For the conservative row, $Ld_1=o(v_1)$.
For the endpoint row, $v_1\ge r$ for large $r$, and
\[
 \frac{Ld_1}{2v_1}
 \le\frac{3A}{2x}+12A2^{-B}<\frac1{100}.
\]
Thus the usual even trace sum over $6x\le\ell\le L$ supplies cycle weight
\begin{equation}\label{eq:endpoint-initial-mass}
 \mu_1:=\sum_{\substack{C\subseteq J_1\\|C|\le L}}2^{-|C|}
 \ge\left(\frac{99}{400}-o(1)\right)\ln(A/6)>20.
\end{equation}
Here $L^22^{-s_1}=o(1)$ controls repeated walks, with $s_j=\lfloor(D_j-1)/2\rfloor$.

At either girth scale, the hypotheses of~\eqref{eq:endpoint-link-estimate} hold for large $r$.
The components meeting $P$ contribute at most $x/z^2$.
Furthermore,
\[
 \frac{2L}{t+1}=(2A+o(1))\frac{x}{z},\quad
 2^{t+1}\le\frac{2x}{z^2},\quad
 578(L/D_j)^2\le578(z/8)^2.
\]
It follows that every link count is at most
\[
 b_r:=3Ax/z.
\]
Since $s_0\ge4Ax/z-2$, the exceptional-pair error satisfies
\[
 \eta_j\le O(L^32^{-Ax/z})=o(1)
 \qquad(j=0,1).
\]
Lemma~\ref{lem:endpoint-two-link-moment} therefore gives
\begin{equation}\label{eq:endpoint-upper-weight}
 \mu_0\le\frac{1+o(1)}{4\varepsilon},
 \qquad W_{2,1}>\mu_1^2-(1+o(1))\mu_1.
\end{equation}

We next obtain a nested chain at the larger girth scale.
For disjoint cycles $C,C'$, call $C'$ far behind $C$ if it lies outside the $P_0$-component of $\Gamma-V(C)$.
This relation is antisymmetric.
A path from $P_0$ to its first vertex on $C$ avoids that far component and hence avoids $C'$.
Connectedness then puts all of $C$ in the $P_0$-component after deleting $C'$.
Thus ordered far pairs have total weight at most $\mu_1^2/2$.
The weight of pairs joined by an edge is at most $\mu_1L^22^{-s_1}=o(\mu_1)$.

Equations~\eqref{eq:endpoint-initial-mass} and~\eqref{eq:endpoint-upper-weight} therefore give a cycle $C$ with partners of total weight at least eight satisfying three conditions.
They lie in the $P_0$-component of $H:=\Gamma-V(C)$, have exactly two links to $C$, and have no edge to $C$.
For each partner, one of the two exterior linking components avoids $P_0$.
It contains a neighbor $u$ of $C$, and the partner separates $u$ from $P_0$ in $H$.
There are at most $L$ choices of $u$.
For one choice, the separating cycles have total weight at least $8/L$.

Take a maximal vertex-disjoint subfamily of these cycles, of size $T$.
Its vertices meet every separating cycle in the family.
The local weight estimate gives
\begin{equation}\label{eq:endpoint-chain-size}
 T\ge\frac{8\,2^{s_1}}{L^3}
 \ge\frac{2^{1-B/2}}{A^3}\frac{\sqrt r}{x^{a/2+3}}.
\end{equation}
The cycles are nested between $u$ and $P_0$ in $H$.
Indeed, two disjoint connected separators cannot cross one another's sides.
Rerouting a $u$--$P_0$ path through the first separator shows that the second lies entirely on one side.
The separation property then gives a consistent order.
Write the ordered cycles as $C_1,\ldots,C_T$, and let $U_i$ be the $u$-component of $H-V(C_i)$.
Then $U_i\subseteq U_{i+1}$ and $C_i\subseteq U_{i+1}$.

Put $q=\lceil16A^2x\rceil+2$.
Choose disjoint intervals $[a',b']$ of indices with $b'-a'=q$, leaving a separator between successive intervals.
For each interval, use the full vertex set
\[
 V(a',b')=U_{b'}\setminus(U_{a'}\cup V(C_{a'})).
\]
It contains all $q-1$ intermediate cycles.
Thus, since $D_1+1\ge x/2$ for large $r$,
\[
 |V(a',b')|\ge(q-1)(D_1+1)\ge8L^2.
\]
Every edge of $H$ leaving this set ends on $C_{a'}$ or $C_{b'}$.
Its boundary in $H$ is at most $2L$.
The chosen full sets are pairwise disjoint.

Discard blocks meeting a neighbor of $C$, a vertex of $S_0$, or a vertex of $P$.
At most $L+|S_0|+c$ blocks are discarded.
For the last term, each connected component of $\Gamma[P]$ can meet at most one block.
It cannot cross a block boundary because every such boundary ends on a separator cycle disjoint from $P$.

Each retained block avoids $P\cup S_0$ and has no neighbor in either set.
An edge to either set would leave the block and end on a separator cycle, which avoids both sets.
The discarded neighbors of $C$ remove the only edges missing from this boundary description in $\Gamma$.
Thus the induced graph on each retained block has girth greater than $D_0$, and degree deficit at most $2L$.
Its order is at most $|V(\Gamma)|\le3r$.
Lemma~\ref{lem:endpoint-block-supply} gives cycle weight at least one in each retained block.

The losses in~\eqref{eq:endpoint-deletion-sizes} are negligible compared with the number of blocks.
Using~\eqref{eq:endpoint-chain-size} and $q+1\le18A^2x$, for large $r$ their number $m$ satisfies
\begin{equation}\label{eq:endpoint-clean-block-count}
 m\ge c_B\frac{\sqrt r}{x^{a/2+4}},
 \qquad c_B:=\frac{2^{-B/2}}{64A^5}.
\end{equation}
Cycles in different blocks are distinct cycles of $J_0$.
Consequently $\mu_0\ge m$.

For the conservative row, this contradicts
\[
 c_0\frac{\sqrt r}{x^{13/2}}
 \le\mu_0\le(1+o(1))\frac{\sqrt r}{4x^7}.
\]
For the endpoint row, choose the absolute constant
\[
 M=\frac1{c_B}=64A^5 2^{B/2}.
\]
Then it contradicts
\[
 c_B\frac{\sqrt r}{x^6}
 \le\mu_0\le(1+o(1))\frac{c_B}{4}\frac{\sqrt r}{x^6}.
\]
This proves both forest estimates.
\end{proof}

\begin{corollary}[Unrestricted endpoint-scale cycle error]
\label{cor:parallel-links-endpoint-cycles}
There is an absolute constant $M$ such that
\[
 \psi(r)\le2^{r-1}+(M+o(1))\frac{2^r(\log r)^6}{\sqrt r}.
\]
In particular, for every fixed $\kappa<1/2$,
\[
 \psi(r)\le2^{r-1}+(1+o(1))2^r r^{-\kappa}.
\]
\end{corollary}

\begin{proof}
Use the direct cubic-extremizer transfer of Proposition~\ref{prop:two-scale-direct-cycle-transfer}.
A shortest cycle gives a connected marked set of size $O(\log r)$, satisfying the new hypotheses.
The cycle-flip injection bounds some cyclic edge count by
\[
 c_e\le2^{r-2}+\frac M2\frac{2^r(\log r)^6}{\sqrt r}.
\]
Deleting that edge gives the same rank recurrence as before.
The geometric tail satisfies
\[
 \sum_{j\le r}2^{j-1}\frac{(\log j)^6}{\sqrt j}
 =(1+o(1))\frac{2^r(\log r)^6}{\sqrt r}.
\]
This proves the first bound.
For each fixed $\kappa<1/2$, one has $(\log r)^6/\sqrt r=o(r^{-\kappa})$.
The second bound follows.
\end{proof}
